\section{Model Cards and Datasheets}

\begin{frame}{Documentation Is Part of the Model}

    A model that ships without documentation forces every downstream user to rediscover its limits
    --- usually in production.

    \vspace{0.8em}

    \begin{columns}[T]
        \begin{column}{0.48\textwidth}
            \begin{tcolorbox}[colback=raiBlue!5, colframe=raiBlue!60, title=Model card]
                \small
                Travels \textbf{with the trained model}. Answers: what is it for, who is it for,
                where does it fail, and how well does it work \emph{per group}.

                \vspace{0.4em}
                Mitchell et al., FAT* 2019.
            \end{tcolorbox}
        \end{column}
        \begin{column}{0.48\textwidth}
            \begin{tcolorbox}[colback=raiGreen!5, colframe=raiGreen!60, title=Datasheet]
                \small
                Travels \textbf{with the dataset}. Answers: why was it collected, from whom, with
                what consent, how was it cleaned, and what should it not be used for.

                \vspace{0.4em}
                Gebru et al., CACM 2021.
            \end{tcolorbox}
        \end{column}
    \end{columns}

    \vspace{1em}

    \begin{block}{Both are borrowed ideas}
        \centering
        \small A model card is a \textbf{datasheet for a component}, as in electronics; a dataset
        datasheet is the same idea one step upstream. Neither is a research artefact --- they are
        engineering hygiene, and increasingly they are what a customer or a partner asks to see
        before they will build on your model.
    \end{block}
\end{frame}

\begin{frame}{What Goes in a Model Card}

    \centering
    \scriptsize
    \begin{tabular}{p{3.2cm} p{8.6cm}}
        \toprule
        \textbf{Section} & \textbf{Content} \\
        \midrule
        Model details      & Owner, date, version, model type, training algorithm, licence, citation, contact. \\
        Intended use       & Primary intended uses and users; \textbf{out-of-scope} uses stated explicitly. \\
        Factors            & Groups, instrumentation, and environments the performance is expected to vary over. \\
        Metrics            & Which measures, which decision threshold, and how uncertainty was estimated. \\
        Evaluation data    & Which datasets, why they were chosen, and how they were preprocessed. \\
        Training data      & Distribution over the same factors, as far as it can be disclosed. \\
        Quantitative analyses & Results \textbf{unitary} (per factor) and \textbf{intersectional} (per combination). \\
        Ethical considerations & Sensitive data, risks to people, mitigations applied, known problematic uses. \\
        Caveats \& recommendations & Untested groups, open concerns, what testing should happen next. \\
        \bottomrule
    \end{tabular}

    \vspace{0.9em}
    \raggedright
    {\scriptsize M.~Mitchell, S.~Wu, A.~Zaldivar, P.~Barnes, L.~Vasserman, B.~Hutchinson, E.~Spitzer,
    I.~D.~Raji and T.~Gebru, ``Model Cards for Model Reporting,'' FAT* 2019, Section 4,
    \href{https://arxiv.org/abs/1810.03993v2}{arXiv:1810.03993v2}.}
\end{frame}

\begin{frame}{A Filled Example (1/2) --- Invented for Teaching}

    \begin{tcolorbox}[colback=gray!4, colframe=raiBlue!60,
                      title={Model card --- \texttt{loan-default-v3}}, fonttitle=\bfseries]
        \scriptsize
        \textbf{Model details.} Gradient-boosted trees (LightGBM, 400 trees, depth 6). Owner:
        Credit Risk team. Version 3.1, trained 2026-05-14 on applications from 2021--2025.
        Contact: \texttt{risk-ml@example.org}.

        \vspace{0.35em}
        \textbf{Intended use.} Produces a default-probability score used to \emph{prioritise manual
        review} of personal-loan applications between \$1k and \$25k in Market X.
        Primary users: underwriting analysts.

        \vspace{0.35em}
        \textbf{Out of scope.} Not for automatic decline without human review. Not for business
        loans, mortgages, or other markets. Not for pricing. Not valid for applicants under 21
        (fewer than 300 in training).

        \vspace{0.35em}
        \textbf{Factors.} Reported over: age band, gender, region, income decile, thin-file vs
        thick-file credit history, and the region\,$\times$\,thin-file intersection.

        \vspace{0.35em}
        \textbf{Metrics.} AUC and recall at the operating threshold $0.31$, chosen to keep the
        manual-review queue at 20\% of volume. 95\% intervals from 1{,}000 bootstrap resamples.

        \vspace{0.35em}
        \textbf{Fairness criterion adopted.} Equal opportunity on the review decision
        (equal TPR across region and thin-file status). Demographic parity was
        \textbf{not} targeted, because base default rates differ materially by region;
        this trade-off was signed off by the Model Risk Committee on 2026-05-02.
    \end{tcolorbox}
\end{frame}

\begin{frame}{A Filled Example (2/2) --- Invented for Teaching}

    \begin{tcolorbox}[colback=gray!4, colframe=raiBlue!60,
                      title={Model card --- \texttt{loan-default-v3} (continued)}, fonttitle=\bfseries]
        \scriptsize
        \textbf{Quantitative analyses.}

        \vspace{0.3em}
        \centering
        \begin{tabular}{lrrrr}
            \toprule
            \textbf{Slice} & \textbf{$n$} & \textbf{AUC} & \textbf{TPR @ 0.31} & \textbf{FPR @ 0.31} \\
            \midrule
            Overall              & 42{,}300 & 0.81 & 0.74 & 0.18 \\
            Region North         & 31{,}500 & 0.82 & 0.76 & 0.17 \\
            Region South         & 10{,}800 & 0.77 & 0.66 & 0.21 \\
            Thin file            &  6{,}100 & 0.71 & 0.58 & 0.24 \\
            \textbf{South $\times$ thin file} &  1{,}450 & \textbf{0.68} & \textbf{0.52} & \textbf{0.27} \\
            \bottomrule
        \end{tabular}

        \vspace{0.5em}
        \raggedright
        \textbf{Ethical considerations.} Region correlates with an attribute protected in Market X.
        Region is not an input, but postcode-derived features are; a proxy audit is run quarterly.
        A wrongly high score delays access to credit, so recall on the worst slice is the primary
        fairness metric.

        \vspace{0.35em}
        \textbf{Caveats and recommendations.} The worst slice is 3.4\% of volume and is
        \textbf{materially worse} than overall --- route it to a specialist queue rather than the
        standard one. Not validated after a policy change to income verification (2026-03).
        Re-audit on 2026-11-14 or after any change to the acquisition channel.
    \end{tcolorbox}
\end{frame}

\begin{frame}[allowframebreaks]{Datasheets for Datasets}

    Seven groups of questions, answered by whoever created the dataset:

    \vspace{0.18em}

    \begin{columns}[T]
        \begin{column}{0.49\textwidth}
            \begin{enumerate}\scriptsize
                \setlength\itemsep{0.1em}
                \item \textbf{Motivation.} Why was it created, by whom, funded by whom?
                \item \textbf{Composition.} What is an instance? How many? What is missing? Does it
                      contain data about people, or confidential/sensitive content?
                \item \textbf{Collection process.} How was it acquired, by whom, over what period?
                      Were the people informed, and did they consent?
                \item \textbf{Preprocessing / cleaning / labelling.} What was done, and is the raw
                      version still available?
            \end{enumerate}
        \end{column}
        \begin{column}{0.49\textwidth}
            \begin{enumerate}\scriptsize
                \setlength\itemsep{0.1em}
                \setcounter{enumi}{4}
                \item \textbf{Uses.} What has it been used for; what should it \textbf{not} be used
                      for?
                \item \textbf{Distribution.} Who gets it, under what licence, with what
                      restrictions?
                \item \textbf{Maintenance.} Who maintains it, how are errors fixed, is there
                      versioning, is there a retention limit?
            \end{enumerate}

            \vspace{0.15em}
            \begin{tcolorbox}[colback=raiAmber!5, colframe=raiAmber!70, boxrule=0.6pt]
                \scriptsize The questions are deliberately uncomfortable. ``We do not know'' is a
                valid answer --- and is itself the finding.
            \end{tcolorbox}

            \vspace{0.12em}
            \begin{tcolorbox}[colback=raiGreen!5, colframe=raiGreen!60, boxrule=0.6pt]
                \scriptsize \textbf{Where these live.} Hugging Face model and dataset cards are
                exactly this format. Keep the card in the \emph{same repository as the training
                code}, generate its quantitative table \emph{from the evaluation script}, and make
                it a merge requirement --- a hand-copied table is stale within a week.
            \end{tcolorbox}
        \end{column}
    \end{columns}

    \vspace{0.12em}
    \begin{block}{Exercise (15 minutes)}
        \centering
        \scriptsize Take the last model you trained. Write its \textbf{intended use}, its
        \textbf{out-of-scope uses}, and one \textbf{quantitative table with slices}.
        Most people discover they cannot fill in the second one.
    \end{block}

    \vspace{0.09em}
    {\scriptsize T.~Gebru, J.~Morgenstern, B.~Vecchione, J.~W.~Vaughan, H.~Wallach, H.~Daum\accent19 e~III
    and K.~Crawford, ``Datasheets for Datasets,'' \emph{Communications of the ACM} 64(12):86--92,
    2021, \href{https://arxiv.org/abs/1803.09010v8}{arXiv:1803.09010v8}.}
\end{frame}

